\documentclass[10pt,a4paper]{amsart}
\usepackage{amsmath,amssymb,amsthm,graphicx}
\usepackage{verbatim}
\newenvironment{reference}{\comment}{\endcomment}
%\newenvironment{reference}{}{}
\newenvironment{slogan}{\comment}{\endcomment}
\newenvironment{history}{\comment}{\endcomment}

% For commutative diagrams we use Xy-pic
\usepackage[all]{xy}

% We use 2cell for 2-commutative diagrams.
\xyoption{2cell}
% Theorem environments.
%
\theoremstyle{plain}
\newtheorem{theorem}[subsection]{Theorem}
\newtheorem{proposition}[subsection]{Proposition}
\newtheorem{lemma}[subsection]{Lemma}

\theoremstyle{definition}
\newtheorem{definition}[subsection]{Definition}
\newtheorem{example}[subsection]{Example}
\newtheorem{exercise}[subsection]{Exercise}
\newtheorem{situation}[subsection]{Situation}

\theoremstyle{remark}
\newtheorem{remark}[subsection]{Remark}
\newtheorem{remarks}[subsection]{Remarks}

\numberwithin{equation}{subsection}

% Macros
%
\def\lim{\mathop{\mathrm{lim}}\nolimits}
\def\colim{\mathop{\mathrm{colim}}\nolimits}
\def\Spec{\mathop{\mathrm{Spec}}}
\def\Hom{\mathop{\mathrm{Hom}}\nolimits}
\def\Ext{\mathop{\mathrm{Ext}}\nolimits}
\def\SheafHom{\mathop{\mathcal{H}\!\mathit{om}}\nolimits}
\def\SheafExt{\mathop{\mathcal{E}\!\mathit{xt}}\nolimits}
\def\Sch{\mathit{Sch}}
\def\Mor{\mathop{\mathrm{Mor}}\nolimits}
\def\Ob{\mathop{\mathrm{Ob}}\nolimits}
\def\Sh{\mathop{\mathit{Sh}}\nolimits}
\def\NL{\mathop{N\!L}\nolimits}
\def\CH{\mathop{\mathrm{CH}}\nolimits}
\def\proetale{{pro\text{-}\acute{e}tale}}
\def\etale{{\acute{e}tale}}
\def\QCoh{\mathit{QCoh}}
\def\Ker{\mathop{\mathrm{Ker}}}
\def\Im{\mathop{\mathrm{Im}}}
\def\Coker{\mathop{\mathrm{Coker}}}
\def\Coim{\mathop{\mathrm{Coim}}}

% Boxtimes
%
\DeclareMathSymbol{\boxtimes}{\mathbin}{AMSa}{"02}

%
% Macros for moduli stacks/spaces
%
\def\QCohstack{\mathcal{QC}\!\mathit{oh}}
\def\Cohstack{\mathcal{C}\!\mathit{oh}}
\def\Spacesstack{\mathcal{S}\!\mathit{paces}}
\def\Quotfunctor{\mathrm{Quot}}
\def\Hilbfunctor{\mathrm{Hilb}}
\def\Curvesstack{\mathcal{C}\!\mathit{urves}}
\def\Polarizedstack{\mathcal{P}\!\mathit{olarized}}
\def\Complexesstack{\mathcal{C}\!\mathit{omplexes}}
% \Pic is the operator that assigns to X its picard group, usage \Pic(X)
% \Picardstack_{X/B} denotes the Picard stack of X over B
% \Picardfunctor_{X/B} denotes the Picard functor of X over B
\def\Pic{\mathop{\mathrm{Pic}}\nolimits}
\def\Picardstack{\mathcal{P}\!\mathit{ic}}
\def\Picardfunctor{\mathrm{Pic}}
\def\Deformationcategory{\mathcal{D}\!\mathit{ef}}
\usepackage{cleveref}
\crefname{originallemma}{original lemma}{original lemmas}
\Crefname{originallemma}{Original lemma}{Original lemmas}
\makeatletter
\newlabel{more-algebra-lemma-dedekind-torsion-free-flat}{{0AUW}{source}{}{sourcepointer-1}{}}
\newlabel{more-algebra-lemma-dedekind-torsion-free-flat@cref}{{[originallemma][0][]0AUW}{[1][0][]source}{}{}{}}
\makeatother
\crefname{originallemma}{original lemma}{original lemmas}
\Crefname{originallemma}{Original lemma}{Original lemmas}
\makeatletter
\newlabel{lemma-flat-pullback-bivariant}{{0B78}{source}{}{sourcepointer-2}{}}
\newlabel{lemma-flat-pullback-bivariant@cref}{{[originallemma][0][]0B78}{[1][0][]source}{}{}{}}
\makeatother
\crefname{originallemma}{original lemma}{original lemmas}
\Crefname{originallemma}{Original lemma}{Original lemmas}
\makeatletter
\newlabel{lemma-push-proper-bivariant}{{0EPK}{source}{}{sourcepointer-3}{}}
\newlabel{lemma-push-proper-bivariant@cref}{{[originallemma][0][]0EPK}{[1][0][]source}{}{}{}}
\makeatother
\crefname{originallemma}{original lemma}{original lemmas}
\Crefname{originallemma}{Original lemma}{Original lemmas}
\makeatletter
\newlabel{lemma-gysin-bivariant}{{0B79}{source}{}{sourcepointer-4}{}}
\newlabel{lemma-gysin-bivariant@cref}{{[originallemma][0][]0B79}{[1][0][]source}{}{}{}}
\makeatother
\title{The exterior product over a base of dimension one is well defined on Chow groups in degree shifted by one}
\author{The Stacks Project Authors}
\date{stacks-project:chow.tex:local-snapshot-20261012; GNU FDL 1.2 or later}
\begin{document}
\maketitle
\noindent\textit{Editorial note.} All delivered mathematical text is the author's own, copied from the pinned Stacks Project chapter file chow.tex; no statement or proof step was rewritten, reordered or removed.\par\medskip
\noindent\textit{Editorial note.} Preamble provenance: the pinned chapter file opens with an input of the shared preamble and carries no preamble text of its own, so the project's single Stacks-wide preamble staged at sources/stacks-chow/preamble.tex is excerpted verbatim as two spans, lines 14-18 (Xy-pic setup) and lines 154-217 (the native theorem declarations and mathematical macros). Nothing was added, removed or redefined; the wrapper class and the amsmath/amssymb/amsthm/graphicx base packages are supplied by the assembler.\par\medskip
\noindent\textit{Editorial note.} The counted statement and proof contain no cross reference, no citation and no included file, so the delivered body is self contained.\par\medskip
\noindent\textit{Editorial note.} Canonical location of the counted proposition: Stacks tag 0FC4 (sources/stacks-chow/pinned-original-source.tex line 14658).\par\medskip
\noindent\textit{Allowed standard prerequisites (editorial).}\begin{itemize}
\item A normal scheme over a Dedekind domain is flat over it, and a point of a Dedekind spectrum is an effective Cartier divisor. More on Algebra, Lemma labelled more-algebra-lemma-dedekind-torsion-free-flat (Stacks tag 0AUW), invoked at proof lines 14686-14687. Supplies that the projection of a component dominating S is flat of the required relative dimension, which is what makes the flat pullback a bivariant class in the first case.
\item Flat pullback, proper pushforward and the Gysin map of an effective Cartier divisor as bivariant classes. Same chapter, Lemmas labelled lemma-flat-pullback-bivariant (source line 6131), lemma-push-proper-bivariant (source line 6167) and lemma-gysin-bivariant (source line 6149), invoked at proof lines 14689-14692 and 14706-14715. Supply the three bivariant operations whose composites c' are capped against classes on Y to realise the cycle level product in each of the two cases.
\end{itemize}\medskip
\ifdefined\hypertarget\hypertarget{sourcepointer-1}{}\fi
\noindent\textit{Source reference.} Original lemma 0AUW: sources/stacks-chow/tags.txt, source lines 13664--13664. Cross-chapter standard prerequisite located in the project tag list staged in the local snapshot; invoked at proof line 14687 of the counted proof. It is retained only as a navigation pointer; every load-bearing step is inside the delivered span.\par\medskip
\ifdefined\hypertarget\hypertarget{sourcepointer-2}{}\fi
\noindent\textit{Source reference.} Original lemma 0B78: sources/stacks-chow/pinned-original-source.tex, source lines 6131--6131. Same-chapter pointer into the author environment lemma at source line 6131; invoked at proof line 14690 of the counted proof. It is retained only as a navigation pointer; every load-bearing step is inside the delivered span.\par\medskip
\ifdefined\hypertarget\hypertarget{sourcepointer-3}{}\fi
\noindent\textit{Source reference.} Original lemma 0EPK: sources/stacks-chow/pinned-original-source.tex, source lines 6167--6167. Same-chapter pointer into the author environment lemma at source line 6167; invoked at proof line 14692 of the counted proof. It is retained only as a navigation pointer; every load-bearing step is inside the delivered span.\par\medskip
\ifdefined\hypertarget\hypertarget{sourcepointer-4}{}\fi
\noindent\textit{Source reference.} Original lemma 0B79: sources/stacks-chow/pinned-original-source.tex, source lines 6149--6149. Same-chapter pointer into the author environment lemma at source line 6149; invoked at proof line 14707 of the counted proof. It is retained only as a navigation pointer; every load-bearing step is inside the delivered span.\par\medskip
\begin{lemma}
\label{lemma-exterior-product-well-defined-dim-1}
The map $\times : \CH_n(X) \otimes_{\mathbf{Z}} \CH_m(Y) \to
\CH_{n + m - 1}(X \times_S Y)$ is well defined.
\end{lemma}

\begin{proof}
Consider $n$ and $m$ cycles $\alpha = \sum_{i \in I} n_i[X_i]$
and $\beta = \sum_{j \in J} m_j[Y_j]$ with $X_i \subset X$ and $Y_j \subset Y$
locally finite families of integral closed subschemes of
$\delta$-dimensions $n$ and $m$. Let $K \subset I \times J$ be the set
of pairs $(i, j) \in I \times J$ such that $X_i$ or $Y_j$ dominates
an irreducible component of $S$.
Then $\{X_i \times_S Y_j\}_{(i, j) \in K}$ is a locally finite
collection of closed subschemes of $X \times_S Y$ of
$\delta$-dimension $n + m - 1$. This means we can indeed consider
$$
\alpha \times \beta =
\sum\nolimits_{(i, j) \in K} n_i m_j [X_i \times_S Y_j]_{n + m - 1}
$$
as a $(n + m - 1)$-cycle on $X \times_S Y$. In this way we obtain an
additive map
$\times : Z_n(X) \otimes_{\mathbf{Z}} Z_m(Y) \to Z_{n + m}(X \times_S Y)$.
The problem is to show that
this procedure is compatible with rational equivalence.

\medskip\noindent
Let $i : X' \to X$ be the inclusion morphism of an integral closed subscheme
of $\delta$-dimension $n$ which dominates an irreducible component
of $S$. Then $p' : X' \to S$ is flat of relative dimension $n - 1$, see
More on Algebra, Lemma \ref{more-algebra-lemma-dedekind-torsion-free-flat}.
Hence flat pullback along $p'$ is an element
$(p')^* \in A^{-n + 1}(X' \to S)$ by
Lemma \ref{lemma-flat-pullback-bivariant}
and hence $c' = i_* \circ (p')^* \in A^{-n + 1}(X \to S)$ by
Lemma \ref{lemma-push-proper-bivariant}.
This produces maps
$$
c' \cap - : \CH_m(Y) \longrightarrow \CH_{m + n - 1}(X \times_S Y)
$$
which sends $[Y']$ to $[X' \times_S Y']_{n + m - 1}$ for any
integral closed subscheme $Y' \subset Y$ of $\delta$-dimension $m$.

\medskip\noindent
Let $i : X' \to X$ be the inclusion morphism of an integral closed subscheme
of $\delta$-dimension $n$ such that the composition $X' \to X \to S$ 
factors through a closed point $s \in S$. Since $s$ is a closed point
of the spectrum of a Dedekind domain, we see that $s$ is an effective
Cartier divisor on $S$ whose normal bundle is trivial. Denote
$c \in A^1(s \to S)$ the Gysin homomorphism, see
Lemma \ref{lemma-gysin-bivariant}. The morphism $p' : X' \to s$
is flat of relative dimension $n$. Hence flat pullback along $p'$
is an element $(p')^* \in A^{-n}(X' \to S)$ by
Lemma \ref{lemma-flat-pullback-bivariant}.
Thus
$$
c' = i_* \circ (p')^* \circ c \in A^{-n}(X \to S)
$$
by Lemma \ref{lemma-push-proper-bivariant}. This produces maps
$$
c' \cap - : \CH_m(Y) \longrightarrow \CH_{m + n - 1}(X \times_S Y)
$$
which for any integral closed subscheme $Y' \subset Y$
of $\delta$-dimension $m$
sends $[Y']$ to either $[X' \times_S Y']_{n + m - 1}$ if $Y'$ dominates
an irreducible component of $S$ or to $0$ if not.

\medskip\noindent
From the previous two paragraphs we conclude
the construction $([X'], [Y']) \mapsto [X' \times_S Y']_{n + m - 1}$
factors through rational equivalence in the second variable, i.e.,
gives a well defined map
$Z_n(X) \otimes_{\mathbf{Z}} \CH_m(Y) \to \CH_{n + m - 1}(X \times_S Y)$.
By symmetry the same is true for the other variable and we conclude.
\end{proof}

\end{document}
